% In this file you should put macros to be used only by
% the web version. Of course they should have a corresponding
% version in macros/print.tex.
% Typically the printed version could have more fancy decorations.
% This will probably be a very short file.

% plasTeX's native \Cref already prints the full capitalized name and links it
% with the number, without equation parentheses. Consume the matching PDF
% settings here: plasTeX does not implement these cleveref configuration commands.
\providecommand{\creflabelformat}[2]{}
\providecommand{\crefname}[3]{}

% plasTeX identifies theorem types by their captions, independently of counters.
\newcommand{\newsharedtheorem}[2]{\newtheorem{#1}[theorem]{#2}}

% Dummy macros that make sense only for web version.
\newcommand{\lean}[1]{}
\newcommand{\discussion}[1]{}
\newcommand{\leanok}{}
\newcommand{\mathlibok}{}
\newcommand{\notready}{}
% % Make sure that arguments of \uses and \proves are real labels, by using invisible refs:
% % latex prints a warning if the label is not defined, but nothing is shown in the pdf file.
% % It uses LaTeX3 programming, this is why we use the expl3 package.
% \ExplSyntaxOn
% \NewDocumentCommand{\uses}{m}
%  {\clist_map_inline:nn{#1}{\vphantom{\ref{##1}}}%
%   \ignorespaces}
% \NewDocumentCommand{\proves}{m}
%  {\clist_map_inline:nn{#1}{\vphantom{\ref{##1}}}%
%   \ignorespaces}
% \ExplSyntaxOff
% Arrow label naming a Lean theorem. MathJax renders the math, so the web version shows the
% name only; the sentence after each such display carries the \ref link.
\newcommand{\leanArrowLabel}[2]{\mathtt{#2}}
